% Part: incompleteness
% Chapter: representability-in-q
% Section: representable-comp

\documentclass[../../../include/open-logic-section]{subfiles}

\begin{document}

\olfileid{inc}{req}{rpc}
\olsection{الدوال القابلة للتمثيل في $\Th{Q}$ قابلة للحساب}

نبدأ بإثبات أن قابلية الدالة للتمثيل في~$\Th{Q}$ تستلزم قابليتها للحساب.
ونقدّم لذلك لمة في الدوال القابلة للتمثيل في~$\Th{Q}$.

\begin{lem}\ollabel{lem:rep-q}
  لتكن $\OLId{latin-ordinary}{f}(\OLId{latin-ordinary}{x}_\OLMathNumeral{0},\dots,\OLId{latin-ordinary}{x}_\OLId{latin-ordinary}{k})$ قابلة للتمثيل في~$\Th{Q}$. حينئذ توجد
  !!a{formula}~$!A_\OLId{latin-ordinary}{f}(\OLId{latin-ordinary}{x}_\OLMathNumeral{0},\dots,\OLId{latin-ordinary}{x}_\OLId{latin-ordinary}{k},\OLId{latin-ordinary}{y})$ تحقق:
  \[
  \Th{Q} \Proves !A_\OLId{latin-ordinary}{f}(\num{\OLId{latin-ordinary}{n}_\OLMathNumeral{0}}, \dots, \num{\OLId{latin-ordinary}{n}_\OLId{latin-ordinary}{k}}, \num{\OLId{latin-ordinary}{m}})
  \quad\text{إذا وفقط إذا}\quad \OLId{latin-ordinary}{m} = \OLId{latin-ordinary}{f}(\OLId{latin-ordinary}{n}_\OLMathNumeral{0}, \dots, \OLId{latin-ordinary}{n}_\OLId{latin-ordinary}{k}).
  \]
\end{lem}

\begin{proof}
أما اتجاه «إذا» فمقتضى
\olref[int]{defn:representable-fn}\olref[int]{defn:rep:a}.
وأما اتجاه «فقط إذا»، فنبرهنه بالخلف. لنفرض
$\Th{Q}\Proves !A_\OLId{latin-ordinary}{f}(\num{\OLId{latin-ordinary}{n}_\OLMathNumeral{0}},\dots,\num{\OLId{latin-ordinary}{n}_\OLId{latin-ordinary}{k}},\num{\OLId{latin-ordinary}{m}})$ مع
$\OLId{latin-ordinary}{m}\neq \OLId{latin-ordinary}{f}(\OLId{latin-ordinary}{n}_\OLMathNumeral{0},\dots,\OLId{latin-ordinary}{n}_\OLId{latin-ordinary}{k})$، ولنضع $\OLId{latin-ordinary}{l}=\OLId{latin-ordinary}{f}(\OLId{latin-ordinary}{n}_\OLMathNumeral{0},\dots,\OLId{latin-ordinary}{n}_\OLId{latin-ordinary}{k})$.
فمن \olref[int]{defn:representable-fn}\olref[int]{defn:rep:a} نحصل على
$\Th{Q}\Proves !A_\OLId{latin-ordinary}{f}(\num{\OLId{latin-ordinary}{n}_\OLMathNumeral{0}},\dots,\num{\OLId{latin-ordinary}{n}_\OLId{latin-ordinary}{k}},\num{\OLId{latin-ordinary}{l}})$.
ومن \olref[int]{defn:representable-fn}\olref[int]{defn:rep:b} نحصل على
$\Th{Q}\Proves
\lforall[\OLId{latin-ordinary}{y}][(!A_\OLId{latin-ordinary}{f}(\num{\OLId{latin-ordinary}{n}_\OLMathNumeral{0}},\dots,\num{\OLId{latin-ordinary}{n}_\OLId{latin-ordinary}{k}},\OLId{latin-ordinary}{y})\lif\num{\OLId{latin-ordinary}{l}}=\OLId{latin-ordinary}{y})]$.
فإذا استعملنا قواعد المنطق وفرضنا
$\Th{Q}\Proves !A_\OLId{latin-ordinary}{f}(\num{\OLId{latin-ordinary}{n}_\OLMathNumeral{0}},\dots,\num{\OLId{latin-ordinary}{n}_\OLId{latin-ordinary}{k}},\num{\OLId{latin-ordinary}{m}})$، لزم
$\Th{Q}\Proves\eq[\num{\OLId{latin-ordinary}{l}}][\num{\OLId{latin-ordinary}{m}}]$.
لكن \olref[bre]{lem:q-proves-neq} يعطينا أيضًا
$\Th{Q}\Proves\eq/[\num{\OLId{latin-ordinary}{l}}][\num{\OLId{latin-ordinary}{m}}]$؛ فتكون $\Th{Q}$ غير متسقة.
وهذا ممتنع: فإن النموذج القياسي يحقق $\Th{Q}$
(انظر \olref[int][def]{def:standard-model})، أي إن $\Sat{\OLId{latin-looped}{N}}{\Th{Q}}$،
وكل نظرية قابلة للتحقيق متسقة بمبرهنة السلامة
(\tagrefs{prfAX/{fol:axd:sou:cor:consistency-soundness},
prfSC/{fol:seq:sou:cor:consistency-soundness},
prfND/{fol:ntd:sou:cor:consistency-soundness},
prfTab/{fol:tab:sou:cor:consistency-soundness}}).
\end{proof}

\begin{lem}
إذا كانت الدالة قابلة للتمثيل في $\Th{Q}$ فهي قابلة للحساب.
\end{lem}

\begin{proof}
نبيّن أولًا وجه القول على سبيل التقريب. فحساب~$\OLId{latin-ordinary}{f}$ يجري بالبحث الآتي:
نسرد كل ما يمكن أن يكون من !!{derivation}s~$\delta$ في لغة الحساب،
وسردها ممكن آليًا. ثم ننظر في كل واحدة: أهي !!a{derivation}
لـ!!a{formula} على الصورة
$!A_\OLId{latin-ordinary}{f}(\num{\OLId{latin-ordinary}{n}_\OLMathNumeral{0}},\dots,\num{\OLId{latin-ordinary}{n}_\OLId{latin-ordinary}{k}},\num{\OLId{latin-ordinary}{m}})$، أعني الـ!!{formula} الممثلة
لـ$\OLId{latin-ordinary}{f}$ في~$\Th{Q}$ بحسب \olref{lem:rep-q}؟ فإن كانت كذلك ثبت
$\OLId{latin-ordinary}{m}=\OLId{latin-ordinary}{f}(\OLId{latin-ordinary}{n}_\OLMathNumeral{0},\dots,\OLId{latin-ordinary}{n}_\OLId{latin-ordinary}{k})$ باللمة \olref{lem:rep-q}، فحصلت قيمة~$\OLId{latin-ordinary}{f}$.
ولا يستمر البحث بلا نهاية؛ إذ
$\Th{Q}\Proves !A_\OLId{latin-ordinary}{f}(\num{\OLId{latin-ordinary}{n}_\OLMathNumeral{0}},\dots,\num{\OLId{latin-ordinary}{n}_\OLId{latin-ordinary}{k}},
\num{\OLId{latin-ordinary}{f}(\OLId{latin-ordinary}{n}_\OLMathNumeral{0},\dots,\OLId{latin-ordinary}{n}_\OLId{latin-ordinary}{k})})$، فلا بد أن يظهر في السرد $\delta$ المطلوبة.

ولم تبلغ هذه الحجة بعدُ غاية الدقة؛ فالإجراء الموصوف يتناول
!!{derivation}s و!!{formula}s، لا الأعداد وحدها، ولم نفصّل وجه إمكان
«سرد جميع الـ!!{derivation}s الممكنة» آليًا. غير أن الطريق إلى ضبطها
قد تقدم: فالحدود و!!{formula}s و!!{derivation}s ترمّز بأعداد غودل،
والدوال العودية العامة تعطينا نموذجًا مضبوطًا للحساب. وقد ثبت أن العلاقة
$\Prf[\Th{Q}](\OLId{latin-ordinary}{d},\OLId{latin-ordinary}{y})$ عودية، بل عودية بدائية؛ ومعناها أن $\OLId{latin-ordinary}{d}$ عدد غودل
لـ!!a{derivation} من بديهيات~$\Th{Q}$ للـ!!{formula} التي عدد غودل لها~$\OLId{latin-ordinary}{y}$.
ونحتاج معها إلى الدالتين العوديتين البدائيتين
$\fn{num}$ (\olref[art][trm]{prop:num-primrec}) و$\fn{Subst}$
(\olref[art][sub]{prop:subst-primrec}).
ومن هذه المعطيات نعرّف~$\OLId{latin-ordinary}{f}$ بالتصغير، فتثبت عودية $\OLId{latin-ordinary}{f}$.

نضع أولًا
\begin{multline*}
  \OLId{latin-looped}{A}(\OLId{latin-ordinary}{n}_\OLMathNumeral{0}, \dots, \OLId{latin-ordinary}{n}_\OLId{latin-ordinary}{k}, \OLId{latin-ordinary}{m}) = \\
  \fn{Subst}(\fn{Subst}(\dots\fn{Subst}(\Gn{!A_\OLId{latin-ordinary}{f}}, \fn{num}(\OLId{latin-ordinary}{n}_\OLMathNumeral{0}), \Gn{\OLId{latin-ordinary}{x}_\OLMathNumeral{0}}),\\ \dots),
  \fn{num}(\OLId{latin-ordinary}{n}_\OLId{latin-ordinary}{k}),  \Gn{\OLId{latin-ordinary}{x}_\OLId{latin-ordinary}{k}}), \fn{num}(\OLId{latin-ordinary}{m}), \Gn{\OLId{latin-ordinary}{y}})
\end{multline*}
وليس وراء تعقيد هذه الكتابة إلا الدالة
$\OLId{latin-looped}{A}(\OLId{latin-ordinary}{n}_\OLMathNumeral{0},\dots,\OLId{latin-ordinary}{n}_\OLId{latin-ordinary}{k},\OLId{latin-ordinary}{m})=
\Gn{!A_\OLId{latin-ordinary}{f}(\num{\OLId{latin-ordinary}{n}_\OLMathNumeral{0}},\dots,\num{\OLId{latin-ordinary}{n}_\OLId{latin-ordinary}{k}},\num{\OLId{latin-ordinary}{m}})}$.

ثم نأخذ العلاقة~$\OLId{latin-looped}{R}(\OLId{latin-ordinary}{n}_\OLMathNumeral{0},\dots,\OLId{latin-ordinary}{n}_\OLId{latin-ordinary}{k},\OLId{latin-ordinary}{s})$، ومعناها أن $(\OLId{latin-ordinary}{s})_\OLMathNumeral{0}$
عدد غودل لـ!!a{derivation} من~$\Th{Q}$ للصيغة
$!A_\OLId{latin-ordinary}{f}(\num{\OLId{latin-ordinary}{n}_\OLMathNumeral{0}},\dots,\num{\OLId{latin-ordinary}{n}_\OLId{latin-ordinary}{k}},\num{(\OLId{latin-ordinary}{s})_\OLMathNumeral{1}})$:
\[
\OLId{latin-looped}{R}(\OLId{latin-ordinary}{n}_\OLMathNumeral{0}, \dots, \OLId{latin-ordinary}{n}_\OLId{latin-ordinary}{k}, \OLId{latin-ordinary}{s}) \quad\text{إذا وفقط إذا}\quad \Prf[\Th{Q}]((\OLId{latin-ordinary}{s})_\OLMathNumeral{0}, \OLId{latin-looped}{A}(\OLId{latin-ordinary}{n}_\OLMathNumeral{0},
\dots, \OLId{latin-ordinary}{n}_\OLId{latin-ordinary}{k}, (\OLId{latin-ordinary}{s})_\OLMathNumeral{1}))
\]
فمتى وجدنا~$\OLId{latin-ordinary}{s}$ تحقق $\OLId{latin-looped}{R}(\OLId{latin-ordinary}{n}_\OLMathNumeral{0},\dots,\OLId{latin-ordinary}{n}_\OLId{latin-ordinary}{k},\OLId{latin-ordinary}{s})$، وجدنا زوج العددين
$(\OLId{latin-ordinary}{s})_\OLMathNumeral{0}$ و$(\OLId{latin-ordinary}{s})_\OLMathNumeral{1}$؛ أولهما، $(\OLId{latin-ordinary}{s})_\OLMathNumeral{0}$، عدد غودل
  لـ!!a{derivation} للصيغة المعبّر عنها اختصارًا بـ$!A_\OLId{latin-ordinary}{f}(\num{\OLId{latin-ordinary}{n}_\OLMathNumeral{0}},\dots,\num{\OLId{latin-ordinary}{n}_\OLId{latin-ordinary}{k}},\num{(\OLId{latin-ordinary}{s})_\OLMathNumeral{1}})$.
فالبحث عن~$\OLId{latin-ordinary}{s}$ هو نظير البحث عن $\OLId{latin-ordinary}{d}$ و$\OLId{latin-ordinary}{m}$ معًا في الحجة غير الصورية.
وهذا البحث عن $\OLId{latin-ordinary}{s}$ تحدده دالة قابلة للحساب بالتّصغير المنتظم.
وانتظام $\OLId{latin-looped}{R}$ بيّن: لكل $\OLId{latin-ordinary}{n}_\OLMathNumeral{0}$، \dots، $\OLId{latin-ordinary}{n}_\OLId{latin-ordinary}{k}$، توجد
!!a{derivation}~$\delta$ تثبت ما تقوله
$\Th{Q}\Proves !A_\OLId{latin-ordinary}{f}(\num{\OLId{latin-ordinary}{n}_\OLMathNumeral{0}},\dots,\num{\OLId{latin-ordinary}{n}_\OLId{latin-ordinary}{k}},
\num{\OLId{latin-ordinary}{f}(\OLId{latin-ordinary}{n}_\OLMathNumeral{0},\dots,\OLId{latin-ordinary}{n}_\OLId{latin-ordinary}{k})})$.
فتصدق $\OLId{latin-looped}{R}(\OLId{latin-ordinary}{n}_\OLMathNumeral{0},\dots,\OLId{latin-ordinary}{n}_\OLId{latin-ordinary}{k},\OLId{latin-ordinary}{s})$ عند
$\OLId{latin-ordinary}{s}=\tuple{\Gn{\delta},\OLId{latin-ordinary}{f}(\OLId{latin-ordinary}{n}_\OLMathNumeral{0},\dots,\OLId{latin-ordinary}{n}_\OLId{latin-ordinary}{k})}$، ونحصل على $\OLId{latin-ordinary}{f}$ بالتعريف:
\[
  \OLId{latin-ordinary}{f}(\OLId{latin-ordinary}{n}_\OLMathNumeral{0},\dots,\OLId{latin-ordinary}{n}_{\OLId{latin-ordinary}{k}}) = (\umin{\OLId{latin-ordinary}{s}}{\OLId{latin-looped}{R}(\OLId{latin-ordinary}{n}_\OLMathNumeral{0}, \dots, \OLId{latin-ordinary}{n}_\OLId{latin-ordinary}{k}, \OLId{latin-ordinary}{s})})_\OLMathNumeral{1}.
  \]
  % تصحيح مصدر (C291-01): وُحّد اسم الصيغة الممثلة وأُثبت واسم الصيغة وغلاف العدد في المثال المختصر.
\end{proof}

\end{document}
